\section{Problema e scelta architetturale}

L'assistente sviluppato per questa tesi consulta servizi con interfacce e formati
diversi. GDACS pubblica feed RSS in XML, il catalogo CDSE usa richieste STAC e
Copernicus EMS restituisce strutture specifiche del servizio. Gestire questi
dettagli direttamente nei nodi del grafo avrebbe legato l'orchestrazione e i
prompt a trasporto, autenticazione, parsing e messaggi di errore delle singole
sorgenti.

Il Model Context Protocol (MCP) è stato adottato per i rami GDACS e Copernicus
EMS allo scopo di introdurre un confine tra il workflow e l'accesso ai dati. La
specifica definisce un protocollo aperto per collegare applicazioni basate su LLM
a strumenti e risorse esterne \cite{mcpSpec}; la proposta originaria insiste in
particolare sulla riusabilità delle integrazioni tra host differenti
\cite{anthropicMCP}.

Il capitolo si concentra sul percorso effettivamente usato nel progetto: host,
client, server e invocazione dei tool.

MCP propone un modello client-server. L'applicazione IA agisce come host: contiene o coordina il modello, presenta l'interfaccia utente e gestisce le policy. Al suo interno usa uno o più client MCP per connettersi a server MCP. Ogni server espone capacità specifiche: strumenti, risorse o prompt. L'agente può scoprire queste capacità e usarle attraverso messaggi standardizzati.

\section{Concetti principali}

\subsection{Host}

L'host coordina modello, client e policy applicative. Nel caso studiato questo
ruolo appartiene al backend FastAPI e al workflow LangGraph; la dashboard nativa
fornisce l'interfaccia con l'utente (Open WebUI resta un client legacy
opzionale).

È l'host a scegliere quali capacità mostrare al modello. Nel progetto la scelta è
ristretta a operazioni di consultazione: i server MCP offrono soltanto letture e
ricerche.

\subsection{Client}

Il client MCP è il componente che mantiene la connessione con un server. In molte architetture è una parte interna dell'host, invisibile all'utente. Il client si occupa di inizializzare la sessione, negoziare capacità, inviare richieste e ricevere risposte.

Se l'host usa più server, mantiene più connessioni. Il backend sviluppato in
questa tesi utilizza un client comune e apre una sessione distinta per GDACS e
Copernicus EMS.

\subsection{Server}

Il server MCP espone funzionalità a cui l'agente può accedere. Può essere locale, per esempio un processo avviato sulla macchina dell'utente, oppure remoto, raggiungibile via rete. Un server può offrire:
\begin{itemize}
    \item \textbf{tools}, cioè azioni invocabili;
    \item \textbf{resources}, cioè dati leggibili e indirizzabili;
    \item \textbf{prompts}, cioè template o procedure testuali riusabili.
\end{itemize}

La distinzione separa dati, azioni e istruzioni riutilizzabili. Nel progetto i
server GDACS e CEMS usano soltanto la prima categoria: espongono come tool le
azioni di lettura e ricerca.

\subsection{Trasporto e messaggi}

MCP usa messaggi JSON-RPC e mantiene il formato logico separato dal trasporto
\cite{mcpSpec}. La specifica prevede, tra gli altri, processi locali collegati
tramite standard input/output e server raggiunti via HTTP. I servizi GDACS e CEMS
del progetto sono container autonomi e comunicano con il backend via HTTP.

\section{Capacità esposte dai server}

\subsection{Tools}

Un tool è una funzione descritta da un nome, da una spiegazione del comportamento
e dallo schema degli argomenti. Nel server GDACS, ad esempio, la ricerca degli
eventi recenti è esposta con un contratto equivalente al seguente:
\begin{verbatim}
tool: list_recent_events
input: { "hazard": "string", "window": "string" }
output: lista normalizzata di eventi GDACS
\end{verbatim}

Descrizione e schema costituiscono le informazioni sulle quali il modello basa
la scelta del tool e dei relativi parametri. Un contratto chiaro si traduce
quindi direttamente in invocazioni più corrette.

Nel prototipo ogni tool corrisponde a un'operazione della sorgente, accetta un
input validato e restituisce errori classificabili. Contratti circoscritti
permettono inoltre di provare parsing e gestione dei guasti separatamente dal
modello linguistico.

\subsection{Resources e prompts}

Una resource rende leggibile un contenuto indirizzabile, di norma identificato da
un URI, mentre i prompts sono template testuali per interazioni ricorrenti. La
prima descrive un dato da recuperare; i secondi forniscono all'host una procedura
o un'istruzione parametrica. Entrambe le capacità fanno parte della specifica; il
flusso realizzato per GDACS e CEMS si basa sui soli tool.

\section{Uso di MCP nel workflow}

Nel progetto i client scoprono i tool all'apertura della sessione. Il ramo
selezionato determina il server e il router specializzato sceglie l'operazione,
così il modello vede solo i tool pertinenti. Dopo la validazione degli argomenti,
il client invia la richiesta e il risultato normalizzato rientra nello stato
LangGraph. La selezione è quindi distribuita fra decisione del modello e vincoli
del grafo.

Quando l'orchestratore seleziona GDACS, il client MCP invia i parametri
normalizzati al tool scelto e riceve una risposta strutturata. Il server si
occupa della richiesta HTTP, del parsing XML e degli errori della sorgente; il
workflow vede soltanto il contratto del tool. Lo stesso schema è usato per le
attivazioni Copernicus EMS. MCP affianca quindi le API remote e ne isola i
dettagli dal resto dell'applicazione.

\section{Confronto e limiti}

\subsection{Integrazione diretta}

Un server MCP, internamente, può chiamare qualunque API: REST, GraphQL, database SQL, SDK proprietari, comandi locali. Verso l'applicazione IA, invece, espone un'interfaccia uniforme.

Rispetto a un'integrazione collocata interamente nel backend, il costo aggiuntivo
è la gestione di sessioni, trasporto e processi separati. In cambio, parsing XML
di GDACS e accesso a CEMS restano confinati nei rispettivi servizi. Il client
STAC, mantenuto direttamente nel backend, offre un termine di confronto concreto:
è più semplice da eseguire, ma rende il backend responsabile dello schema e degli
errori dell'API remota.

\subsection{Tool calling}

Il tool calling tradizionale descrive le funzioni che una singola applicazione
rende disponibili al modello. MCP opera tra l'host e server separati, dai quali
quelle funzioni possono essere scoperte. L'host può quindi presentare al modello
un tool ottenuto da un server MCP e tradurre l'invocazione nella corrispondente
richiesta di protocollo. I due meccanismi intervengono in punti diversi della
stessa catena e sono complementari.

\subsection{Costi introdotti}

Ogni sorgente conserva la propria semantica: eventi GDACS e attivazioni CEMS
hanno campi, tempi di aggiornamento e condizioni di errore diversi. MCP
introduce inoltre due servizi da distribuire e monitorare. Per un prototipo con
poche funzioni tale costo è visibile; il vantaggio sta nel poter modificare una
sorgente lasciando parsing e dipendenze fuori dal workflow.

\section{Sicurezza e gestione operativa}

\subsection{Autenticazione e autorizzazione}

Avvio dei server, accesso ai dati, gestione delle credenziali e revoca dei
permessi sono responsabilità dell'host e dell'infrastruttura, più che del
protocollo di integrazione.

I server sono raggiungibili sulla rete interna dei container e ricevono
configurazione e credenziali dall'ambiente. Il server stabilisce le operazioni
disponibili; il workflow limita ulteriormente quelle invocabili nel ramo
corrente. Poiché tutti i tool sono di sola lettura, il prototipo può fare a meno
di un flusso di conferma per azioni irreversibili.

\subsection{Prompt injection e dati esterni}

Feed e pagine web possono contenere testo che contraddice le istruzioni
dell'applicazione. I contenuti recuperati vengono trattati come dati, mentre
piano, tool disponibili, elenco degli indirizzi GDACS ammessi e schemi di output
sono controllati dal codice. Questa separazione riduce la superficie d'attacco,
pur senza escludere ogni forma di prompt injection.

\subsection{Osservabilità}

Il client registra servizio, durata, esito e codice diagnostico delle chiamate,
omettendo credenziali e prompt completi. MCP offre un punto uniforme nel quale
raccogliere queste informazioni, mentre log e metriche sono implementati
dall'applicazione.

\section{Criteri adottati nel progetto}

I tool implementati sono piccoli, di sola lettura e dotati di schemi di input
espliciti. Le credenziali e gli indirizzi interni dei servizi restano nella
configurazione applicativa. I risultati includono invece URL pubblici delle
risorse, identificativi e timestamp utili alla verifica. I test dei client
coprono risposte vuote, payload malformati, timeout ed errori HTTP temporanei.
Queste proprietà derivano dalla progettazione dei singoli server e dell'host,
più che dall'adozione di MCP in sé.

Nel sistema proposto MCP ha un compito delimitato: separa i servizi GDACS
e Copernicus EMS dal backend agentico e offre un'interfaccia uniforme per
scoprirne e invocarne i tool. Il catalogo satellitare STAC, usato per
acquisizioni e heatmap nel workspace, resta un'integrazione diretta, così come
l'accesso della Discovery al catalogo NebulaGraph. La convivenza dei due
approcci mostra in concreto quando MCP conviene come scelta architetturale.
